\documentclass[10pt,a4paper]{amsart}
\usepackage[margin=19mm]{geometry}
\usepackage{amsmath,amssymb,mathtools}
\usepackage[scr=rsfs,cal=euler]{mathalfa}
\usepackage{xfrac}
\usepackage[unicode,hidelinks,bookmarks=false]{hyperref}
\newcommand{\ie}{i.e.\ }
\newcommand{\cf}{cf.\ }
\newcommand{\resp}{respectively\ }
\newcommand{\Subsection}{\subsection}
\providecommand{\nobreakdash}{\nobreak\hbox{-}}
\setlength{\emergencystretch}{2em}
\multlinegap0pt
\usepackage{amssymb}
\newtheorem{prop}{Proposition}
\newtheorem{lemma}{Lemma}
\newcommand{\dd}{\underline{d}}
\renewcommand{\geq}{\geqslant}
\renewcommand{\leq}{\leqslant}
\title{Topological K-theory of quasi-BPS categories of symmetric quivers with potential}
\author{Tudor P\u{a}durariu, Yukinobu Toda}
\date{Source: arXiv:2309.08432v3 (CC BY 4.0)}
\begin{document}
\maketitle

\noindent\emph{Source and licence.} Topological K-theory of quasi-BPS categories of symmetric quivers with potential. Version arXiv:2309.08432v3 (CC BY 4.0), distributed by arXiv under the Creative Commons Attribution 4.0 International licence, http://creativecommons.org/licenses/by/4.0/, as recorded in the arXiv licence metadata for this exact version and shown on the abstract page https://arxiv.org/abs/2309.08432v3. Full licence notice: \texttt{licenses/arxiv-2309.08432-CC-BY-4.0.txt}.\par\smallskip

\noindent\textit{Editorial scope.} Counted unit: the article's prop prop:delta (source0.tex lines 5294--5297). The article's complete original proof of it is delivered with it (source0.tex lines 5298--5344, one \texttt{proof} environment of 47 lines). No mathematics is written, repaired or re-derived: the statement and the proof are the article's own lines, in the article's own notation, and no retained line is substituted or re-flowed.  Retained uncounted as the source-local context the proof needs: lines 5185--5194; lines 5212--5221.  Nothing was omitted from any retained span, so the package carries no omission record.  No retained line contains a citation, so the wrapper carries no bibliography and no bibliography entry.  No retained line contains a figure environment, an image-inclusion command or drawing code, so no image is delivered and the package declares no asset dependency.  Editorial formatting only, in addition: an AMS \texttt{amsart} preamble that restates the article's own declarations and macros used by the retained lines, namely the environments \texttt{prop}, \texttt{lemma} on separate counters, together with the standard packages those lines need.  The retained environments print in this document as Proposition 1, Lemma 1 in order of appearance, whereas the article prints the same items with the numbering of its own full document.  The complete native source of the article as distributed in the arXiv e-print for this exact version is bundled unchanged as \texttt{sources/147085/source0.tex}.
\begin{prop}\label{computationsetsdv}
    Let $Q=(I,E)$ be a quiver with an even number of edges between 
    any two different vertices and an odd number of loops at every vertex. 
    Let $(d,v)\in \mathbb{N}^I\times\mathbb{Z}$. 
    The set $S^d_v$ consists of partitions $\mathbf{d}=(d_i)_{i=1}^k$ such that
    \begin{align*}
        v \cdot \frac{\dd_i}{\dd} \in \mathbb{Z} \ \mbox{ for all } 1\leq i\leq k. 
    \end{align*}
    In particular, if $\gcd(v,\dd)=1$, then $S^d_v$ contains only the one term partition of $d$.
\end{prop}

\begin{lemma}\label{lem:onevertex}
Let $Q$ be a quiver with one vertex and $2e$-loops for some $e \geq 1$. 
Then $S_v^d$ consists of partitions $\mathbf{d}=(d_i)_{i=1}^k$ such that, for all $1\leq i\leq k$, we have 
\begin{align*}
    \frac{1}{2}d_i\left( \sum_{j<i}d_j-\sum_{j>i}d_j \right)+\frac{vd_i}{d} \in \mathbb{Z}. 
\end{align*}
Moreover,
$S_{v}^d=\{d\}$ if and only if $\gcd(d, v)=1$ and $d \not\equiv 2 \pmod 4$, 
    or if $\gcd(d, v)=2$ and $d \equiv 2 \pmod 4$. 
\end{lemma}

\begin{prop}\label{prop:delta}
For any symmetric quiver $Q=(I, E)$, there is $\delta \in M(d)_{\mathbb{R}}^{W_d}$ such that 
$S_{\delta}^d=\{d\}$. 
\end{prop}

\begin{proof}
    The condition $\{d\} \in S_{\delta}^d$ is equivalent to that 
    $\langle 1_d, \delta \rangle \in \mathbb{Z}$. 
    In the case that $\lvert I \rvert=1$, then $\delta=v\tau_d$ for some $v \in \mathbb{Z}$. 
    By Proposition~\ref{computationsetsdv} and Lemma~\ref{lem:onevertex}, we have 
    $S_{v}^d=\{d\}$ for either $v=1$ or $v=2$. 
    
    In the case that $\lvert I \rvert \geq 2$, 
    we first take $\delta \in M(d)_{\mathbb{R}}^{W_d}$
    such that $v:=\langle 1_d, \delta \rangle \in \mathbb{Z}$, and set $\delta'=\delta-v\tau_d$. 
    Then $\langle 1_d, \delta' \rangle =0$. 
Note that the space of 
    $\delta' \in M(d)_{\mathbb{R}}^{W_d}$ satisfying $\langle 1_d, \delta' \rangle=0$
    is $\mathbb{R}^{\lvert I \rvert-1}(:=M(d)_{\mathbb{R}, 0}^{W_d})$, hence uncountable. 
Therefore for a very general choice of $\delta'$, 
we have $n_{\lambda}/2+\langle \lambda, \delta' \rangle \notin \mathbb{Z}$
for any $\lambda$ such that the hyperplane 
\begin{align*}H_{\lambda}:=\{x\in M(d)_{\mathbb{R}, 0}^{W_d}: \langle \lambda, x \rangle=0\} 
\subset M(d)_{\mathbb{R}, 0}^{W_d}=\mathbb{R}^{\lvert I \rvert-1}
\end{align*}
is of codimension one. We have 
$H_{\lambda}=\mathbb{R}^{\lvert I \rvert-1}$ if and only 
if the associated partition 
$d=d_1+\cdots+d_k$ is proportional to $d$, i.e. 
by writing $d=m d_0$ with $d_0$ a primitive dimension 
vector and $m \in \mathbb{Z}$, we have $d_i=m_i d_0$
where $m=m_1+\cdots+m_k$ is a partition of $m$. 
Then we have 
\begin{align}\label{compute:lambda}
    \frac{n_{\lambda}}{2}+\langle \lambda, \delta \rangle 
    =\frac{1}{2}\left\langle \nu, (R(d_0)-\mathfrak{g}(d_0))
    \otimes \mathrm{End}(W)^{\nu >0} \right\rangle +\langle \nu, v\tau_m
    \rangle. 
\end{align}
Here $W$ is a $m$-dimensional vector space 
and $\nu$ is an antidominant cocharacter of the maximal 
torus of $GL(W)$ with associated partition $(m_i)_{i=1}^k$. 
Therefore, we are reduced to the same computation
for the one-vertex case
with $e$-loops, where 
$e=\dim R(d_0)-\dim \mathfrak{g}(d_0)+1$,
and 
(\ref{compute:lambda}) is not an integer
for any such $\lambda$ for $v=1$ if $e$ is odd, for $v=1$ if $e$ is even and  
$m \not\equiv 2 \pmod 4$, for $v=2$ if 
$e$ is even and $m\equiv 2 \pmod 4$. 
\end{proof}

\end{document}
